\section*{Conclusion}
\addcontentsline{toc}{section}{Conclusion}

The order-four arithmetic and the finite-group operator algebra are not
the obstruction.  The edge-projected Green form is now uniformly
dominated for every $N$ with $K_G=\pi^2/4$, and it converges to the
chiral Green form in local $H^2$ quadratic-form norm.  Exact compatible
dyadic edge-mode isometries now give the charged graph-defect rate
$N^{-2}$ for fields and adjoints.  The full
crossed-product/expectation limit is proved, and the holomorphic $c=8$
endpoint is rigid once a local conformal limit with the charged field
exists.  The unresolved point is the support-preserving identification
of this edge-mode limit with the full microscopic CAR local scaling
algebra. For the original microscopic collar, v1026 now proves the
relaxed uniform TEL-b norm below $2.995906<3$, including CF, DG and
finite-grid aliases, without a finite-plateau inference. Relative locality and
outerness are proved.  The remaining in-house algebraic obligation is
the Buchholz--Verch upper inclusion with the binding FE-GEN/gauge-code
qualification. The conditional sector/actual-word frame theorem in v1030
is a concrete sufficient interface, not a proof that the microscopic
one-edge algebra satisfies it. Joint-adjoint control and high-output tails
remain indispensable.

\begin{thebibliography}{99}

\bibitem{BDL}
D.~Buchholz, C.~D'Antoni and R.~Longo,
\emph{Nuclear maps and modular structures II: applications to quantum
field theory},
Commun.\ Math.\ Phys.\ \textbf{129} (1990), 115--138.

\bibitem{BuchholzWichmann}
D.~Buchholz and E.~H.~Wichmann,
\emph{Causal independence and the energy-level density of states in
local quantum field theory},
Commun.\ Math.\ Phys.\ \textbf{106} (1986), 321--344.

\bibitem{BuchholzVerch}
D.~Buchholz and R.~Verch,
\emph{Scaling algebras and renormalization group in algebraic quantum
field theory},
Rev.\ Math.\ Phys.\ \textbf{7} (1995), 1195--1239.

\bibitem{CKLW}
S.~Carpi, Y.~Kawahigashi, R.~Longo and M.~Weiner,
\emph{From vertex operator algebras to conformal nets and back},
Mem.\ Amer.\ Math.\ Soc.\ \textbf{254} (2018), no.~1213,
\href{https://arxiv.org/abs/1503.01260}{arXiv:1503.01260}.

\bibitem{CarpiTomassini}
S.~Carpi and T.~Tomassini,
\emph{Energy bounds for vertex operator algebra extensions},
Lett.\ Math.\ Phys.\ \textbf{113} (2023), 87,
\href{https://arxiv.org/abs/2303.14097}{arXiv:2303.14097}.

\bibitem{DongMason}
C.~Dong and G.~Mason,
\emph{Holomorphic vertex operator algebras of small central charge},
Pacific J.\ Math.\ \textbf{213} (2004), 253--266,
\href{https://arxiv.org/abs/math/0203005}{arXiv:math/0203005}.

\bibitem{Gui}
B.~Gui,
\emph{Energy bounds condition for intertwining operators of types
$B$, $C$, and $G_2$ unitary affine vertex operator algebras},
Trans.\ Amer.\ Math.\ Soc.\ \textbf{372} (2019), 7371--7424,
\href{https://arxiv.org/abs/1809.07003}{arXiv:1809.07003}.

\bibitem{KLM}
Y.~Kawahigashi, R.~Longo and M.~M\"uger,
\emph{Multi-interval subfactors and modularity of representations in
conformal field theory},
Commun.\ Math.\ Phys.\ \textbf{219} (2001), 631--669.

\bibitem{LongoRehren}
R.~Longo and K.-H.~Rehren,
\emph{Nets of subfactors},
Rev.\ Math.\ Phys.\ \textbf{7} (1995), 567--597.

\bibitem{OsborneStottmeister}
T.~J.~Osborne and A.~Stottmeister,
\emph{Conformal field theory from lattice fermions},
Commun.\ Math.\ Phys.\ \textbf{398} (2023), 219--283,
\href{https://arxiv.org/abs/2107.13834}{arXiv:2107.13834}.

\bibitem{Peller}
V.~V.~Peller,
\emph{Hankel operators of class $\mathfrak S_p$ and their applications
(rational approximation, Gaussian processes, the problem of majorizing
operators)},
Math.\ USSR Sbornik \textbf{41} (1982), 443--479.

\bibitem{ToledanoLaredo}
V.~Toledano Laredo,
\emph{Fusion of positive energy representations of $L\operatorname{Spin}_{2n}$},
Ph.D.\ thesis, University of Cambridge, 1997; analytic energy-bound
results as used in the type-$D$ loop-group literature.

\bibitem{Wassermann}
A.~Wassermann,
\emph{Operator algebras and conformal field theory III: fusion of
positive energy representations of $LSU(N)$ using bounded operators},
Invent.\ Math.\ \textbf{133} (1998), 467--538.

\end{thebibliography}

\end{document}
